\documentclass[10pt,a4paper]{amsart}
\usepackage[margin=19mm]{geometry}
\usepackage{amsmath,amssymb,mathtools}
\usepackage[scr=rsfs,cal=euler]{mathalfa}
\usepackage{xfrac}
\usepackage[unicode,hidelinks,bookmarks=false]{hyperref}
\newcommand{\ie}{i.e.\ }
\newcommand{\cf}{cf.\ }
\newcommand{\resp}{respectively\ }
\newcommand{\Subsection}{\subsection}
\providecommand{\nobreakdash}{\nobreak\hbox{-}}
\setlength{\emergencystretch}{2em}
\multlinegap0pt
\usepackage{amssymb}
\usepackage{mathtools}
\newtheorem{lemma}{Lemma}
\title{Supersaturation for Eventown via Generator Switching}
\author{Zicheng Han, Xiande Zhang, Yuhao Zhao}
\date{Source: arXiv:2608.16321v1 (CC BY 4.0)}
\begin{document}
\maketitle

\noindent\emph{Source and licence.} Supersaturation for Eventown via Generator Switching. Version arXiv:2608.16321v1 (CC BY 4.0), distributed by arXiv under the Creative Commons Attribution 4.0 International licence, http://creativecommons.org/licenses/by/4.0/, as recorded in the arXiv licence metadata for this exact version and shown on the abstract page https://arxiv.org/abs/2608.16321v1. Full licence notice: \texttt{licenses/arxiv-2608.16321-CC-BY-4.0.txt}.\par\smallskip

\noindent\textit{Editorial scope.} Counted unit: the article's lemma lem:switching-fourier (source0.tex lines 343--364). The article's complete original proof of it is delivered with it (source0.tex lines 366--521, one \texttt{proof} environment of 156 lines). No mathematics is written, repaired or re-derived: the statement and the proof are the article's own lines, in the article's own notation, and no retained line is substituted or re-flowed.  Retained uncounted as the source-local context the proof needs: lines 275--288.  Nothing was omitted from any retained span, so the package carries no omission record.  No retained line contains a citation, so the wrapper carries no bibliography and no bibliography entry.  No retained line contains a figure environment, an image-inclusion command or drawing code, so no image is delivered and the package declares no asset dependency.  Editorial formatting only, in addition: an AMS \texttt{amsart} preamble that restates the article's own declarations and macros used by the retained lines, namely the environments \texttt{lemma} on separate counters, together with the standard packages those lines need.  The retained environments print in this document as Lemma 1 in order of appearance, whereas the article prints the same items with the numbering of its own full document.  The complete native source of the article as distributed in the arXiv e-print for this exact version is bundled unchanged as \texttt{sources/207822/source0.tex}.
\begin{lemma}[Maximum eventown subfamilies]\label{lem:maximum-eventown}
Every eventown subfamily of \(\mathcal V\) is contained in a generator.  More
precisely, let \(\mathcal R\subseteq\mathcal V\), let \(\mathcal E\) be a
maximum eventown subfamily of \(\mathcal R\), and let \(\mathcal W\) be any
generator containing \(\mathcal E\).  Then
\[
 \mathcal R\cap\mathcal W=\mathcal E,
\]
and \(\mathcal W\) maximizes \(|\mathcal R\cap\mathcal W'|\) over all
generators \(\mathcal W'\).  Conversely, if \(\mathcal W\) is a generator
maximizing \(|\mathcal R\cap\mathcal W|\), then
\(\mathcal R\cap\mathcal W\) is a maximum eventown subfamily of
\(\mathcal R\).
\end{lemma}

\begin{lemma}[Switching--Fourier estimate]\label{lem:switching-fourier}
Let \(\mathcal R\subseteq\mathcal V\) be a nonempty subset. Let \(\mathcal E\) be a
maximum eventown subfamily of \(\mathcal R\), and write \(\mathcal Y=\mathcal R\setminus\mathcal E\). 
Then
\begin{equation}
 \bigl(|\mathcal E||\mathcal Y|-2e(\mathcal E,\mathcal Y)\bigr)^2
 \le
 2^{\lfloor n/2\rfloor+1}|\mathcal E|e(\mathcal E,\mathcal Y).
\label{eq:switching-fourier}
\end{equation}
Consequently, we have
\begin{equation}
 |\mathcal Y|
 \le
 \frac{2e(\mathcal E,\mathcal Y)}{|\mathcal E|}
 +
 \sqrt{
       \frac{2^{\lfloor n/2\rfloor+1}e(\mathcal E,\mathcal Y)}{|\mathcal E|}
      }.
\label{eq:remainder-bound}
\end{equation}
\end{lemma}

\begin{proof}
By Lemma~\ref{lem:maximum-eventown}, choose a generator \(\mathcal W\)
containing \(\mathcal E\).  Then
\(\mathcal R\cap\mathcal W=\mathcal E, \mathcal Y=\mathcal R\setminus\mathcal W\), 
and \(\mathcal W\) maximizes \(|\mathcal R\cap\mathcal W'|\) over all
generators \(\mathcal W'\).

Partition \(\mathcal Y\) by the cosets of \(\mathcal W\) in
\(\mathcal V\) other than \(\mathcal W\) itself.  For an occupied coset
\(\Gamma=x+\mathcal W\), define
\[
 \lambda_\Gamma:\mathcal W\longrightarrow\mathbb F_2,
 \qquad
 \lambda_\Gamma(w)=\langle x,w\rangle.
\]
This definition is independent of the representative \(x\). Indeed, replacing
\(x\) by \(x+w_0\), with \(w_0\in\mathcal W\), does not change the value
because \(\mathcal W\) is totally isotropic.  The functional is nonzero,
for otherwise \(x\in\mathcal W^\perp\cap\mathcal V=\mathcal W\), contrary
to \(\Gamma\ne\mathcal W\).  Moreover, distinct cosets induce distinct functionals:
if \(x+\mathcal W\) and \(x'+\mathcal W\) induce the same functional, then
\(x-x'\in\mathcal W^\perp\cap\mathcal V=\mathcal W\), so the cosets are
equal.

Set \(K_\Gamma=\ker\lambda_\Gamma,
 m_\Gamma=|\mathcal Y\cap\Gamma|\), and \(b_\Gamma =\bigl|\{y\in\mathcal E:\lambda_\Gamma(y)=1\}\bigr|\).
Since \(\lambda_\Gamma\) is nonzero, \(K_\Gamma\) is a hyperplane in
\(\mathcal W\), and \(\Gamma\) is the disjoint union of two affine
\(K_\Gamma\)-cosets.  We call these two cosets the affine halves of
\(\Gamma\).

Every vector in \(\Gamma\) has exactly \(b_\Gamma\) neighbors in
\(\mathcal E\).  Indeed, if \(z,z'\in\Gamma\), then
\(z'-z\in\mathcal W\), and for every \(y\in\mathcal E\subseteq\mathcal W\),
\[
 \langle z',y\rangle
 =
 \langle z,y\rangle+\langle z'-z,y\rangle
 =
 \langle z,y\rangle.
\]
It follows that
\begin{equation}
 e(\mathcal E,\mathcal Y)=\sum_\Gamma m_\Gamma b_\Gamma.
\label{eq:q-cosets}
\end{equation}

We next obtain the switching bound that controls the weights
\(m_\Gamma\).  Let \(H\) be one of the two affine halves of an occupied
coset \(\Gamma\).  If \(\mathcal Y\cap H=\varnothing\), there is nothing to
prove.  Otherwise choose \(z\in\mathcal Y\cap H\).  For every
\(w\in\mathcal W\), the equality
\(\langle z,w\rangle=\lambda_\Gamma(w)\) holds, and hence
\[
 \mathcal K:=\mathcal W\cap z^\perp=K_\Gamma.
\]
Define \(\mathcal L=\mathcal K\oplus\langle z\rangle\). The space \(\mathcal K\) is totally isotropic, \(z\) is orthogonal to
\(\mathcal K\), and \(\langle z,z\rangle=0\) because
\(z\in\mathcal V\).  Also \(z\notin\mathcal W\), so
\(\dim\mathcal L=\dim\mathcal W\).  Thus \(\mathcal L\) is a generator.
Moreover, we have \(\mathcal L\cap\mathcal W=\mathcal K\) and
\[
 \mathcal L\setminus\mathcal W=z+\mathcal K=H.
\]
Since \(\mathcal W\) maximizes \(|\mathcal R\cap\mathcal W'|\) over all
generators \(\mathcal W'\), we have
\[
 |\mathcal R\cap\mathcal L|
 \le
 |\mathcal R\cap\mathcal W|
 =|\mathcal E|.
\]
Using the displayed decomposition of \(\mathcal L\), we have the exact
identity
\[
 \mathcal R\cap\mathcal L
 =
 (\mathcal E\cap\mathcal K)
 \sqcup
 (\mathcal Y\cap H).
\]
Since \(|\mathcal E\cap\mathcal K|=|\mathcal E|-b_\Gamma\), it follows
that
\[
 |\mathcal Y\cap H|\le b_\Gamma.
\]
The two affine halves cover \(\Gamma\), so \(m_\Gamma\le 2b_\Gamma\) and hence
\begin{equation}
 \sum_\Gamma m_\Gamma^2
 \le
 2\sum_\Gamma m_\Gamma b_\Gamma
 =2e(\mathcal E,\mathcal Y).
\label{eq:weight-energy}
\end{equation}

We now turn to the Fourier step.  The quantity on the left side of
\eqref{eq:switching-fourier} can be written as a weighted sum of Fourier
coefficients.  Let \(\chi_{\mathcal E}\) denote the indicator function of
\(\mathcal E\) on \(\mathcal W\).  For \(\lambda\in\mathcal W^*\), use the
unnormalized coefficient
\[
 \widehat{\chi_{\mathcal E}}(\lambda)
 =
 \sum_{w\in\mathcal W}
 \chi_{\mathcal E}(w)(-1)^{\lambda(w)}.
\]
For an occupied coset \(\Gamma\),
\[
 \widehat{\chi_{\mathcal E}}(\lambda_\Gamma)
 =
 |\mathcal E|-2b_\Gamma.
\]
Parseval's identity in this normalization gives
\[
 \sum_{\lambda\in\mathcal W^*}
 \widehat{\chi_{\mathcal E}}(\lambda)^2
 =
 |\mathcal W|
 \sum_{w\in\mathcal W}\chi_{\mathcal E}(w)^2
 =
 2^{\lfloor n/2\rfloor}|\mathcal E|.
\]
Since distinct occupied cosets give distinct functionals, we deduce
\begin{equation}
 \sum_\Gamma (|\mathcal E|-2b_\Gamma)^2
 \le
 2^{\lfloor n/2\rfloor}|\mathcal E|.
\label{eq:fourier-energy}
\end{equation}
Finally, by \eqref{eq:q-cosets},
\[
 |\mathcal E||\mathcal Y|-2e(\mathcal E,\mathcal Y)
 =
 \sum_\Gamma m_\Gamma(|\mathcal E|-2b_\Gamma).
\]
The Cauchy--Schwarz inequality, followed by \eqref{eq:weight-energy} and
\eqref{eq:fourier-energy}, yields
\[
 \bigl(|\mathcal E||\mathcal Y|-2e(\mathcal E,\mathcal Y)\bigr)^2
 \le
 \left(\sum_\Gamma m_\Gamma^2\right)
 \left(\sum_\Gamma(|\mathcal E|-2b_\Gamma)^2\right)
 \le
 2^{\lfloor n/2\rfloor+1}|\mathcal E|e(\mathcal E,\mathcal Y).
\]
This proves \eqref{eq:switching-fourier}.  Taking square roots gives
\[
 |\mathcal E||\mathcal Y|-2e(\mathcal E,\mathcal Y)
 \le
 \bigl||\mathcal E||\mathcal Y|-2e(\mathcal E,\mathcal Y)\bigr|
 \le
 \sqrt{2^{\lfloor n/2\rfloor+1}|\mathcal E|e(\mathcal E,\mathcal Y)}.
\]
Adding \(2e(\mathcal E,\mathcal Y)\) and dividing by \(|\mathcal E|\) gives
\eqref{eq:remainder-bound}.
\end{proof}

\end{document}
